\section{Conexión nonádica, extensión del calendario y memoria de transporte}
\label{cap:calendario}

La cronología del TPK combina una fase observable, un registro dodecafásico
y una memoria de las transiciones realizadas. Su relación se determina por
la composición de los desplazamientos. El paso por la frontera de fase
incrementa un cociente; la conservación de ese cociente distingue los
retornos sucesivos. Esta estructura proporciona la base aritmética de la
holonomía nonádica y de su representación bilateral.

Conviene fijar desde el comienzo los tres objetos que intervienen. La fase
pertenece a un ciclo de nueve posiciones. El calendario finito retiene doce
vueltas de ese ciclo. La historia enriquecida conserva, además, el número
entero de vueltas y los registros de posición, orientación, hoja, acarreo y
frontera. Cada reducción conserva una observación determinada de la misma
evolución; la información disponible se precisa mediante sus proyecciones.

\subsection{Incidencia de los modos aritméticos}
\label{sec:cal-incidencia}

\subsubsection{Selección, permanencia y sucesión cíclica}

Sean $\Sigma$ y $\Pi$ los modos correspondientes a las evaluaciones aditiva
y multiplicativa de APP. El selector de fase recorre, en cada bloque de tres
eventos, la sucesión $(+1,-1,0)$. Su acción sobre el modo $m$ es
\[
 +1:m\longmapsto\Sigma,\qquad
 -1:m\longmapsto\Pi,\qquad
 0:m\longmapsto m.
\]
El evento neutral conserva el modo precedente. Partiendo del cierre cíclico
del bloque, los tres modos registrados son, por tanto,
\[
 (\Sigma,\Pi,\Pi).
\]
La tercera posición conserva la etiqueta de modo multiplicativo
seleccionada en el evento anterior. En el emisor, ese evento neutral
incrementa únicamente el contador $Z$. Esta distinción permite
construir la incidencia directamente desde las transiciones del calendario.

\begin{proposicion}[Incidencia modal del bloque trítico]
\label{prop:cal-incidencia}
Al contar los pares consecutivos del bloque, incluido su cierre, y ordenar
los modos como $(\Sigma,\Pi)$, se obtiene
\begin{equation}
 F_{\rm av}=\begin{pmatrix}0&1\\1&1\end{pmatrix}.
 \label{eq:cal-incidencia}
\end{equation}
En una sucesión de $3q$ eventos compuesta por $q$ repeticiones del bloque,
el recuento de transiciones es $qF_{\rm av}$.
\end{proposicion}
\begin{proof}
Los pares ordenados son $\Sigma\to\Pi$, $\Pi\to\Pi$ y
$\Pi\to\Sigma$. Cada uno aparece una vez; $\Sigma\to\Sigma$
aparece cero veces. Éstas son, respectivamente, las entradas superior
derecha, inferior derecha, inferior izquierda y superior izquierda de la
matriz. Cada repetición añade exactamente los mismos tres pares, incluido
el enlace desde el último modo de un bloque hasta el primero del siguiente.
La suma de los recuentos es $qF_{\rm av}$.
\end{proof}

Las longitudes $9$, $27$ y $108$ producen así $3F_{\rm av}$,
$9F_{\rm av}$ y $36F_{\rm av}$. Estas matrices cuentan eventos de un
itinerario periódico ya fijado. Las potencias de $F_{\rm av}$ cuentan
concatenaciones admitidas por su incidencia. Se trata de dos operaciones
matemáticas distintas sobre la misma matriz: suma de ocurrencias y
composición de transiciones.

\subsubsection{Concatenación de incidencias y recurrencia inducida}

Definamos enteros $a_0=0$, $a_1=1$ y
$a_{n+1}=a_n+a_{n-1}$ para $n\geq1$. Su aparición se obtiene de la
multiplicación de la matriz ya construida.

\begin{proposicion}[Potencias de la incidencia modal]
\label{prop:cal-potencias}
Para todo $n\geq1$,
\[
 F_{\rm av}^{\,n}=
 \begin{pmatrix}a_{n-1}&a_n\\a_n&a_{n+1}\end{pmatrix},
 \qquad
 a_{n-1}a_{n+1}-a_n^2=(-1)^n.
\]
\end{proposicion}
\begin{proof}
Para $n=1$, la fórmula reproduce \eqref{eq:cal-incidencia} porque
$a_2=1$. Si la fórmula vale para $n$, la multiplicación por $F_{\rm av}$
da
\[
 \begin{pmatrix}
 a_n&a_{n-1}+a_n\\a_{n+1}&a_n+a_{n+1}
 \end{pmatrix}
 =\begin{pmatrix}a_n&a_{n+1}\\a_{n+1}&a_{n+2}\end{pmatrix}.
\]
La inducción prueba la primera identidad. Tomando determinantes y usando
$\det F_{\rm av}=-1$ se obtiene la segunda.
\end{proof}

En particular, $F_{\rm av}^{2}=F_{\rm av}+I$. Esta identidad permite
normalizar un generador hiperbólico mediante
\begin{equation}
 J_{-1}=\frac{2F_{\rm av}^{2}-3I}{\sqrt5}
       =\frac{2F_{\rm av}-I}{\sqrt5},
 \qquad J_{-1}^{2}=I.
 \label{eq:cal-generador-hiperbolico}
\end{equation}
En efecto, $(2F_{\rm av}-I)^2=4(F_{\rm av}+I)-4F_{\rm av}+I=5I$.
El coeficiente $5$ procede aquí de la identidad matricial. La representación
cuadrática y la extensión de Euler utilizan este generador después de su
construcción modal.

La recurrencia $(a_n)$ admite posteriormente su identificación con la
sucesión de Fibonacci. El orden demostrativo de este apartado es explícito:
selector trítico, sucesión de modos, conteo de aristas, incidencia y
recurrencia. La identificación de la sucesión conserva ese orden. La
generación regional coinductiva tiene además sus propios emisores, filtros,
elevaciones y condiciones de supervivencia; la recurrencia modal precisa
una de sus realizaciones estructurales.

\subsection{Extensión aritmética del calendario}

\subsubsection{Descomposición en fase y posición dodecafásica}

Escribimos $C_q=\ZZ/q\ZZ$ para el grupo aditivo de residuos módulo $q$.
La posición de un calendario de $108=12\cdot9$ eventos posee una
descomposición única
\[
 t=9b+r,\qquad 0\leq b<12,\quad0\leq r<9,
\]
para el representante $0\leq t<108$. La coordenada $r$ es la fase
nonádica; $b$ identifica la posición de la vuelta en el registro
dodecafásico. La escala $108$ reúne esos dos registros según la ley de
composición con acarreo.

\begin{teorema}[Extensión del calendario con acarreo]
\label{thm:cal-extension}
Las aplicaciones
\[
 \iota:C_{12}\longrightarrow C_{108},\quad[b]\longmapsto[9b],
 \qquad p:C_{108}\longrightarrow C_9,\quad[t]\longmapsto[t]_9
\]
determinan una sucesión exacta
\begin{equation}
 0\longrightarrow C_{12}\xrightarrow{\ \iota\ }C_{108}
 \xrightarrow{\ p\ }C_9\longrightarrow0.
 \label{eq:cal-extension}
\end{equation}
En las coordenadas $(b,r)$, la adición viene dada por
\begin{equation}
 (b,r)\oplus(b',r')=
 \left(\left[b+b'+\left\lfloor\frac{r+r'}9\right\rfloor\right]_{12},
 [r+r']_9\right).
 \label{eq:cal-ley-grupo}
\end{equation}
Esta extensión carece de sección homomorfa.
\end{teorema}
\begin{proof}
La igualdad $[9b]=0$ en $C_{108}$ equivale a $12\mid b$, por lo que
$\iota$ es inyectiva. La aplicación $p$ es sobreyectiva y su núcleo
consiste exactamente en las clases múltiplos de nueve, imagen de
$\iota$. Esto demuestra la exactitud.

Para sumar dos representantes, se escribe
\[
 9b+r+9b'+r'
 =9\left(b+b'+\left\lfloor\frac{r+r'}9\right\rfloor\right)
   +[r+r']_9.
\]
La reducción del primer coeficiente módulo doce da
\eqref{eq:cal-ley-grupo}. Supongamos ahora que existe una sección
homomorfa $s:C_9\to C_{108}$. El representante $a$ de $s([1])$
satisface $a\equiv1\pmod9$ y $9a\equiv0\pmod{108}$. La segunda
condición implica $12\mid a$, de donde $a\equiv0\pmod3$; la primera
implica $a\equiv1\pmod3$. La contradicción demuestra la última
afirmación.
\end{proof}

\subsubsection{Cociclo, inversa y restitución}

El término
\[
 \kappa_9(r,s)=\left\lfloor\frac{r+s}9\right\rfloor,
 \qquad 0\leq r,s<9,
\]
registra el cruce del umbral. Su identidad de composición es
\begin{equation}
 \kappa_9(r,s)+\kappa_9([r+s]_9,u)
 =\kappa_9(s,u)+\kappa_9(r,[s+u]_9).
 \label{eq:cal-cociclo}
\end{equation}
Ambos lados son $\lfloor(r+s+u)/9\rfloor$: basta sustituir
$r+s=9\kappa_9(r,s)+[r+s]_9$, o la descomposición análoga de $s+u$.
La identidad asegura que el cociente acumulado es independiente de la
parentización de tres sumandos.

La inversa de $(b,r)$ queda determinada por
\[
 -(b,r)=
 \begin{cases}
 ([-b]_{12},0),&r=0,\\
 ([-b-1]_{12},9-r),&1\leq r<9.
 \end{cases}
\]
El término adicional $-1$ representa la corrección exacta de frontera al
invertir un residuo interior. Por ejemplo, el representante $t=17$ tiene
coordenadas $(1,8)$. Su inversa es $(-2,1)$ en las coordenadas con
primer componente módulo doce: $9(-2)+1=-17$. El residuo invertido y
el cociente corregido restituyen conjuntamente el entero opuesto.

\subsection{Elevación entera y holonomía nonádica}

\subsubsection{Fase finita y contador de vueltas}

La división euclídea se extiende a todo $k\in\ZZ$:
\[
 k=9m+r,\qquad m=\left\lfloor\frac k9\right\rfloor,
 \quad0\leq r<9.
\]
El par $(m,r)$ determina exactamente $k$. Su reducción
$(m,r)\mapsto([m]_{12},r)$ recupera el calendario finito. En la
elevación entera, una vuelta de nueve pasos actúa por
\begin{equation}
 (r,m)\longmapsto(r,m+1).
 \label{eq:cal-vuelta-entera}
\end{equation}
Doce vueltas restituyen las dos coordenadas finitas y trasladan el
contador entero en doce unidades. La memoria enriquecida incluye ese
contador junto con el registro ordenado de las transiciones.

\subsubsection{Composición de prolongaciones admisibles}

En el nivel $k$, representamos el estado regional enriquecido por
\[
 \widehat x_k=(B_k,C_k,p_k,c_k,\Sigma_k,W_k,g_k,m_k).
\]
$B_k$ reúne los bloques regionales; $C_k$ es el cilindro compatible;
$p_k$ es el prefijo; $c_k$ registra el acarreo; $\Sigma_k$ conserva la
frontera y $W_k$ la incidencia. Las últimas coordenadas son la fase y
el contador de vueltas. Los operadores de prolongación actúan sobre
este estado completo, con las dependencias de cada componente.

\begin{definicion}[Holonomía nonádica del sistema de prolongación]
Sean $\mathcal R_k\subseteq\widehat X_k\times\widehat X_{k+1}$
las relaciones admisibles producidas por el TPK. La composición de una
vuelta es
\[
 \Gamma_{9,k}=\mathcal R_{k+8}\circ\cdots\circ\mathcal R_k.
\]
Para $s\geq1$, la composición de $s$ vueltas es
\[
 \Gamma_{9s,k}=
 \Gamma_{9,k+9(s-1)}\circ\cdots\circ\Gamma_{9,k}.
\]
\end{definicion}

La composición relacional retiene las prolongaciones que satisfacen
simultáneamente las restricciones del estado. Una trayectoria admisible
individualiza una sucesión de tales transiciones. La notación distingue
la relación de admisibilidad, la trayectoria elegida y la observación
finita de su fase.

\begin{proposicion}[Retorno de fase y avance de memoria]
\label{prop:cal-memoria}
En una trayectoria cuya regla de vuelta es
$g_{k+9}=g_k$, $m_{k+9}=m_k+1$, toda composición admisible de $s$
vueltas satisface
\[
 g_{k+9s}=g_k,\qquad m_{k+9s}=m_k+s,
 \qquad [m_{k+9s}]_{12}=[m_k+s]_{12}.
\]
\end{proposicion}
\begin{proof}
El caso $s=1$ es la regla de vuelta. Al componer una vuelta adicional
con la trayectoria ya construida, la fase se conserva y el contador
recibe otra unidad. La inducción demuestra las dos primeras igualdades;
la proyección módulo doce da la tercera.
\end{proof}

El registro permite reconocer cada retorno como una transformación de
la historia. La fase observable expresa una periodicidad, mientras que
el contador y las coordenadas de frontera distinguen sus realizaciones
sucesivas. Las demostraciones de prolongación regional especifican
adicionalmente qué cilindros sobreviven y qué prefijos quedan determinados.

\subsection{Representación bilateral del transporte de fase}

\subsubsection{Desplazamiento unitario y grupo de vueltas}

La elevación del índice admite una representación en
$\mathcal H_{\rm ind}=\ell^2(\ZZ)$. Su base ortonormal se escribe
$e_k=e_{r,m}$ cuando $k=9m+r$. Definimos $Te_k=e_{k+1}$, es decir,
\[
 Te_{r,m}=\begin{cases}
 e_{r+1,m},&r<8,\\e_{0,m+1},&r=8.
 \end{cases}
\]
La transformación permuta biyectivamente una base ortonormal y se
extiende por continuidad a un operador unitario. La inversa lleva
$e_{r,m}$ a $e_{r-1,m}$ cuando $r>0$ y lleva $e_{0,m}$ a
$e_{8,m-1}$.

\begin{teorema}[Representación fiel del contador de vueltas]
\label{thm:cal-representacion}
Sea $S=T^9$. Entonces
\[
 S^se_{r,m}=e_{r,m+s},\qquad s\in\ZZ.
\]
La aplicación $s\mapsto S^s$ es una representación unitaria fiel de
$(\ZZ,+)$.
\end{teorema}
\begin{proof}
Nueve pasos conservan el residuo y aumentan el cociente en una unidad.
La misma identidad aplicada al inverso proporciona la fórmula para
exponentes negativos. Las potencias satisfacen
$S^{s+t}=S^sS^t$. Si $S^s=I$, la igualdad
$e_{r,m+s}=e_{r,m}$ fuerza $s=0$ por la independencia de los vectores
de la base. Esto prueba la fidelidad.
\end{proof}

El espacio de índices es una realización del contador. El espacio de
estados del TPK incorpora además los datos aritméticos e incidenciales.
La inversión unitaria del índice se aplica a su recubrimiento bilateral;
la inversión de una transición completa utiliza el archivo de esa
transición y sus condiciones de admisibilidad.

\subsubsection{Operador de memoria y corrección de frontera}

Definamos, sobre los vectores de soporte finito,
\[
 D_Me_k=\left\lfloor\frac k9\right\rfloor e_k,
 \qquad \mathcal Ie_k=e_{-k},
 \qquad P_0e_k=\begin{cases}e_k,&9\mid k,\\0,&9\nmid k.\end{cases}
\]
Este dominio es denso y estable bajo las operaciones consideradas.

\begin{proposicion}[Identidades de memoria e inversión]
En dicho dominio se cumplen
\begin{equation}
 [D_M,S]=S,\qquad
 \mathcal IT\mathcal I=T^{-1},\qquad
 \mathcal ID_M\mathcal I=-D_M-(I-P_0).
 \label{eq:cal-inversion}
\end{equation}
\end{proposicion}
\begin{proof}
Si $k=9m+r$, entonces
$D_MS e_k=(m+1)e_{k+9}$ y $SD_Me_k=me_{k+9}$, lo que prueba
el conmutador. La composición $\mathcal IT\mathcal I$ lleva
$e_k$ a $e_{k-1}$. Finalmente,
\[
 \left\lfloor\frac{-k}9\right\rfloor
 =\begin{cases}-m,&r=0,\\-m-1,&r>0.\end{cases}
\]
La acción diagonal del último miembro de \eqref{eq:cal-inversion}
coincide con esa expresión en cada vector de la base. La linealidad
concluye la demostración.
\end{proof}

La corrección $I-P_0$ localiza la contribución de frontera en las fases
interiores. Al invertir el recorrido, el cociente y el residuo se
transforman de manera conjunta. Esa identidad expresa aritméticamente
la información que desaparece cuando se observa únicamente la fase.

\subsection{Capacidad de refinamiento y vacancias aritméticas}

\subsubsection{Comparación exacta de bloques ternarios y decimales}

Un bloque de seis trits dispone de $3^6=729$ configuraciones en su
espacio ambiente. Un bloque decimal de tres cifras dispone de $10^3=1000$
posiciones. Sus cardinales definen una capacidad aritmética de
codificación. Para $n\geq0$, sea
\begin{equation}
 \mathcal C_n=\max\{a\in\NN:1000^a\leq729^{n+1}\}.
 \label{eq:cal-capacidad-entera}
\end{equation}
La definición compara enteros y puede evaluarse exactamente. En una
representación logarítmica posterior,
\[
 \rho=\log_{1000}729=\log_{10}9,\qquad
 \delta=1-\rho,\qquad
 \mathcal C_n=\lfloor(n+1)\rho\rfloor.
\]
La razón $\rho$ pertenece a $(0,1)$ y es irracional. Si
$\rho=p/q$ con enteros positivos, se tendría $10^p=9^q$;
la factorización prima de ambos miembros lo impide.

\begin{definicion}[Registro de vacancias de capacidad]
Se definen
\[
 Z_n=\lfloor(n+1)\delta\rfloor,\qquad
 z_n=\lfloor(n+1)\delta\rfloor-\lfloor n\delta\rfloor,
 \quad n\geq0.
\]
El valor $z_n\in\{0,1\}$ registra la retención de capacidad en el
paso correspondiente. Se tiene $z_0=0$ en el origen elegido.
\end{definicion}

\begin{teorema}[Balance de capacidad y vacancias]
\label{thm:cal-capacidad}
Para $n\geq0$ en la primera identidad y $n\geq1$ en la segunda,
\begin{equation}
 \mathcal C_n+Z_n=n,
 \qquad \mathcal C_n-\mathcal C_{n-1}=1-z_n.
 \label{eq:cal-capacidad}
\end{equation}
En particular, para $0\leq a<b$,
\[
 \mathcal C_b-\mathcal C_a+\sum_{n=a+1}^{b}z_n=b-a.
\]
\end{teorema}
\begin{proof}
La irracionalidad de $\delta$ implica
$\lceil(n+1)\delta\rceil=\lfloor(n+1)\delta\rfloor+1$.
Por tanto,
\[
 \lfloor(n+1)(1-\delta)\rfloor
 =n+1-\lceil(n+1)\delta\rceil=n-Z_n.
\]
La diferencia entre dos igualdades consecutivas da la segunda identidad
para $n\geq1$. La suma telescópica demuestra la fórmula en un intervalo.
\end{proof}

El registro $Z_n$ puede obtenerse mediante $n-\mathcal C_n$,
utilizando sólo la comparación entera
\eqref{eq:cal-capacidad-entera}. Las funciones logarítmicas expresan
la misma construcción en coordenadas analíticas. La primera retención
de capacidad ocurre entre $n=20$ y $n=21$:
\[
 \mathcal C_{19}=19,\quad\mathcal C_{20}=20,\quad
 \mathcal C_{21}=20,\quad\mathcal C_{22}=21.
\]
La historia recibe un bloque adicional mientras la capacidad decimal
permanece constante. La vacancia registra precisamente ese desfase.

\subsubsection{Diferenciación de los índices operativos}

La fase del emisor finito, la fase de prolongación regional y la capacidad
decimal pertenecen a índices declarados. La conexión nonádica conserva
la fase tras nueve pasos de su propio índice; la vacancia conserva una
capacidad durante un paso del índice de bloques. La comparación entre
ambas se expresa por
\[
 [\mathcal C_n]_q=\bigl[[n]_q-[Z_n]_q\bigr]_q,
 \qquad q=9\ \text{o}\ 108.
\]
El desfase depende del registro acumulado de vacancias. Conservarlo
permite transformar una lectura de fase en la otra sin identificar
indebidamente los dos calendarios.

La estructura obtenida enlaza selección modal, composición del acarreo,
holonomía y refinamiento. La memoria cumple funciones precisas en cada
nivel: restituye cocientes, distingue vueltas, registra transiciones y
controla la capacidad de codificación. Estas operaciones constituyen el
soporte temporal y archivístico sobre el que actúan los lectores regionales
y las realizaciones de conservación de la unidad.
